\section{Threats to Validity}\label{sec:threats_to_validity}

\noindent\textbf{Construct validity.}
MR coverage provides a relation-level adequacy signal over an inferred MR space rather than a complete behavioral specification. Because this space is inferred from public artifacts, it may omit intended behaviors, include imperfect relations, or depend on the granularity at which behavioral relations are represented. For example, a broad MR may be decomposed into multiple finer-grained relations, leading to different \textit{Touch} and \textit{Cover} values. We mitigate these risks through taxonomy-constrained inference, multi-artifact grounding, hybrid deterministic--semantic alignment, and manual validation of sampled MRs and alignment decisions. The oracle-strengthening study introduces an additional threat because injected assertions only approximate inferred relations; we reduce this risk through relation-specific templates and manual inspection of failing executions. MR coverage should therefore be interpreted with respect to the inferred MR space rather than as an absolute measure of behavioral completeness.

\vspace{4pt}
\noindent\textbf{Internal validity.}
MR inference and semantic alignment rely partly on LLM reasoning and may be affected by model-specific biases. We mitigate this threat through fixed prompts, fixed inference parameters, a fixed taxonomy, identical analysis procedures, and evaluation across multiple LLMs. Although the inferred MR sets differ across models, the main findings remain consistent across the evaluated configurations.

\vspace{4pt}
\noindent\textbf{External validity.}
Our evaluation covers four mature open-source UI component libraries spanning the React and Vue ecosystems. The findings may not generalize to proprietary frameworks or application-specific UI components. In addition, the oracle-strengthening study was conducted on two libraries with reproducible execution environments, and different libraries or assertion templates may lead to different quantitative results. The findings should therefore be interpreted within the studied setting, although the proposed assessment framework is not tied to a specific UI library or framework.